\begin{theorem}[High stationary classes: wild quadratic case]
\label{mod:parameters-highquadratic}
Let $K/F$ be prime cyclic of degree two and residue degree one.  Fix a
positive lower break $t$, a generating integral uniformizer as in the
ramification results, and write
\[
 T=t+1,\qquad m=a(\chi_F)\geq T,\qquad v=m-T,
 \qquad r=\left\lfloor\frac T2\right\rfloor,
 \qquad s=\left\lceil\frac T2\right\rceil.
\]
Assume that $\chi_F$ is minimal in its whole norm-character orbit, and put
\[
 \chi_K=\chi_F\circ N_{K/F},\qquad
 \psi_K=\psi_F\circ\operatorname{Tr}_{K/F}.
\]
For the actual stationary conductor decompositions write
\[
 m=2d+\varepsilon,qquad
 m_K=a(\chi_K)=2d_K+\varepsilon_K,qquad
 q_F=d+\varepsilon,qquad q_K=d_K+\varepsilon_K.
\]
Then the actual upstairs conductor and all of its stationary depths are
\[
 \boxed{\ m_K+T=2m\ },\qquad
 d_K=v+r,\qquad \varepsilon_K=T\bmod 2,
 \qquad q_K=v+s=m-r.
\]
In particular,
\[
 q_K-v=s,\quad q_F\geq s,\quad d\leq d_K,\quad
 r\leq t,\quad r+s=T,\quad
 \varepsilon=m\bmod2.
\]
Thus the parity of the upstairs critical line is the parity of $T$, whereas
the downstairs and twisted critical lines have the parity of $m$.
The positive wild quadratic break also forces
\(\operatorname{char}(k_F)=\operatorname{char}(k_K)=2\).

Let $\tau\ne1$ be a norm character.  Its exact conductor is $T$, with
stationary decomposition $(r,T\bmod2)$.  The actual datum for
$\tau\chi_F$ has conductor exactly $m$ and hence retains the downstairs
decomposition $(d,\varepsilon)$; this includes the boundary $m=T$.
Choose $\gamma_F$ with
\[
 v_F(\gamma_F)=m+n_F(\psi_F).
\]
The ordered common-denominator pair is
\[
 (\gamma_F,\gamma_K)=
 (\gamma_F,\operatorname{algebraMap}_{F,K}(\gamma_F)),
 \qquad
 v_K(\gamma_K)=m_K+n_K(\psi_K).
\]
Lean constructs the norm-polynomial precision at the three exact pairs
$(q_K,q_F)$, $(m_K,m)$, and $(d_K,d)$, including when $d>t$; no norm
representative at precision $d$ is used.

The manuscript's field representatives are supplied by one simultaneous
existential package.  First choose a unit $\beta$ representing
\(\operatorname{St}(\chi_F,\psi_F;\gamma_F)\) in
\(\mathcal O_F/\mathfrak p_F^d\).  The break-level norm theorem gives
independent choices
\[
 \beta_1\in\mathcal O_K^\times,
 \qquad \delta\in U_F^t,
 \qquad \boxed{\ \beta=\delta N_{K/F}(\beta_1)\ }
\]
as an exact identity.  The unit correction $\delta$ remains in every
subsequent formula; it is never replaced by $1$.

For the stationary numerator of $\tau$ at ambient conductor $m$, choose
its class in
\[
 \mathfrak p_F^v/\mathfrak p_F^{v+r}.
\]
Only at the permitted subcritical precision $r\leq t$, the norm
representative lemma yields $\alpha_1\in K^\times$ with
\[
 v_K(\alpha_1)=v,qquad
 \alpha=N_{K/F}(\alpha_1),qquad
 c=\alpha\delta,qquad
 [c]=\operatorname{St}_m(\tau,\psi_F;\gamma_F)
 \quad\text{in }
 \mathfrak p_F^v/\mathfrak p_F^{v+r}.
\]
Equivalently, relative to
\(\gamma_\tau=\gamma_F/c\), whose order is $T+n_F(\psi_F)$, the stationary
coefficient class of $\tau$ is $[1]\in\mathcal O_F/\mathfrak p_F^r$.
No representative in any of these quotients is declared canonical.

Put
\[
 u=\frac{\beta_1}{\alpha_1}\in K^\times,
 \qquad n=\frac{\beta}{c}\in F^\times.
\]
The exact ratio and valuation identities are
\[
 N_{K/F}(u)=n,qquad cn=\beta,qquad
 v_K(u)=v_F(n)=T-m=-v.
\]
Define only for this existential package
\[
 z=\frac{\beta_1\,c}{\alpha_1}=c\,u,\qquad
 \beta_K^0=\beta-z=c(n-u),qquad
 \beta_{\tau}^0=\beta+c=c(n+1),
\]
where scalars in the formula for $\beta_K^0$ are mapped to $K$.
These field elements are selected representatives, not canonical lifts.

For every $y\in\mathfrak p_K^s$, the exact quadratic identity is
\[
 N_{K/F}(1+y)=1+\operatorname{Tr}_{K/F}(y)+N_{K/F}(y),
 \qquad
 \operatorname{Tr}_{K/F}(y),N_{K/F}(y)\in\mathfrak p_F^s.
\]
Since $\tau$ is trivial on norms, its stationary linearization gives the
sign-sensitive conversion
\[
 \psi_F\!\left(\frac{cN(y)}{\gamma_F}\right)
 =\psi_F\!\left(-\frac{c\operatorname{Tr}(y)}{\gamma_F}\right).
\]
Consequently the compatible stationary classes are precisely
\[
 \boxed{\quad
 [\beta_K^0]
   =\operatorname{St}(\chi_K,\psi_K;\gamma_K)
   =P^*_{K/F;\gamma_F,\gamma_K}
       ([\beta])
 \quad}
 \quad\text{in }\mathcal O_K/\mathfrak p_K^{d_K},
\]
and
\[
 \boxed{\quad
 [\beta_{\tau}^0]
   =\operatorname{St}(\tau\chi_F,\psi_F;\gamma_F)
 \quad}
 \quad\text{in }\mathcal O_F/\mathfrak p_F^d.
\]
Here $P^*$ is the adjoint of the truncated norm for the displayed ordered
denominator pair.  Both boxes are quotient-class equalities; neither is an
equality between arbitrary representatives in $K$.

At the boundary $m=T$, minimal-orbit stationary noncancellation proves
\[
 \beta+c\notin\mathfrak p_F,
\]
and hence $1+n\ne0$.  For $m>T$ the same denominator statement follows
from $v_F(n)=-v<0$.  Thus no boundary conclusion is obtained merely by
pretending that the two stationary representatives have unique lifts.

Finally $v_K(z)=v$, and the trace-depth theorem gives
\(\operatorname{Tr}_{K/F}(z)\in\mathfrak p_F^d\), while
\(N_{K/F}(z)=\beta c\) exactly.  Retaining the trace sign in
\[
 N_{K/F}(\beta-z)
  =\beta^2-\beta\operatorname{Tr}_{K/F}(z)+N_{K/F}(z)
\]
therefore yields exactly the product-class congruence
\[
 [N_{K/F}(\beta_K^0)]=[\beta(\beta+c)]
 \quad\text{in }\mathcal O_F/\mathfrak p_F^d.
\]
The characteristic-two interface makes no choice of a refinement.  For
every explicitly supplied normalized refinement $q$, it retains the
positive polar convention
\[
 q(x+y)=q(x)q(y)\psi_0(xy)
 \]
and the Frobenius normalization $q(x^2)=q(x)$ required downstream.

\LeanFile{LanglandsFirstMainLemma/Parameters/HighQuadratic.lean}
\lean{LanglandsFirstMainLemma.wildQuadraticHigh_depth_relations}
\lean{LanglandsFirstMainLemma.wildQuadratic_normPolynomialValue_eq_trace_add_norm}
\lean{LanglandsFirstMainLemma.wildQuadratic_norm_sub}
\lean{LanglandsFirstMainLemma.highParameter_wildQuadratic_conductor_relation}
\lean{LanglandsFirstMainLemma.highParameter_wildQuadratic_residueCharacteristic}
\lean{LanglandsFirstMainLemma.highParameter_wildQuadratic_residueCharacteristic_upstairs}
\lean{LanglandsFirstMainLemma.highParameter_wildQuadratic_commonDenominator}
\lean{LanglandsFirstMainLemma.highParameter_wildQuadratic_precision}
\lean{LanglandsFirstMainLemma.highParameter_wildQuadratic_trace_norm_mem}
\lean{LanglandsFirstMainLemma.wildQuadraticHighTauDecomposition}
\lean{LanglandsFirstMainLemma.highParameter_wildQuadratic_twist_conductor}
\lean{LanglandsFirstMainLemma.WildQuadraticHighRepresentatives}
\lean{LanglandsFirstMainLemma.highParameter_wildQuadratic_exists_representatives}
\lean{LanglandsFirstMainLemma.WildQuadraticHighRepresentatives.norm_u_eq_n}
\lean{LanglandsFirstMainLemma.WildQuadraticHighRepresentatives.u_order}
\lean{LanglandsFirstMainLemma.WildQuadraticHighRepresentatives.n_order}
\lean{LanglandsFirstMainLemma.WildQuadraticHighRepresentatives.tau_unit_class}
\lean{LanglandsFirstMainLemma.WildQuadraticHighRepresentatives.conversion}
\lean{LanglandsFirstMainLemma.WildQuadraticHighRepresentatives.twist_class}
\lean{LanglandsFirstMainLemma.WildQuadraticHighRepresentatives.upstairs_class}
\lean{LanglandsFirstMainLemma.WildQuadraticHighRepresentatives.betaK_adjoint_class}
\lean{LanglandsFirstMainLemma.WildQuadraticHighRepresentatives.boundary_noncancellation}
\lean{LanglandsFirstMainLemma.WildQuadraticHighRepresentatives.norm_betaK_class}
\lean{LanglandsFirstMainLemma.WildQuadraticHighParameterData}
\lean{LanglandsFirstMainLemma.highParameter_wildQuadratic}
\lean{LanglandsFirstMainLemma.highParameter_wildQuadratic_charTwoNormalization}
\leanok
\SourceText{Propositions prop:high-parameter-table and prop:wild-quadratic-high-parameters.}
\uses{mod:parameters-stationaryclassundernorm,
mod:parameters-minimalorbitstationary,
mod:ramification-normrepresentatives,
mod:finitefield-chartworefinement}
\FormalizationContract
The principal theorem returns a nonempty simultaneous representative
package together with the actual conductor relation, exact norm-polynomial
precision, common ordered denominator pair, downstairs twist class,
upstairs stationary class and denominator-sensitive adjoint class, product
quotient, exact correction norm, trace depth, denominator nonvanishing, and
the minimal-orbit boundary statement.  The only exact norm representative
is selected at $r=\lfloor(t+1)/2\rfloor\leq t$; the unrelated unit
$\beta_1$ is obtained at the break with the indispensable factor
$\delta\in U_F^t$.  All stationary equalities live in their stated lattice
quotients.  The result neither selects a canonical stationary lift nor a
canonical characteristic-two refinement, and it imports no later
high-product, high-table, low-parameter, phase-reduction, or case theorem.
\end{theorem}
